Многие задачи механики сплошной среды и физики приводят к исследованию дифференциальных уравнений с частными производными второго порядка (ДУЧП2П). Так, например.

1) При изучении различных видов волн (упругих, звуковых, электромагнитных) мы приходим к \textsl{волновому уравнению}:
\[
\frac{\partial^2 u}{\partial t^2} = c^2 \left( \frac{\partial^2 u}{\partial x^2} + \frac{\partial^2 u}{\partial y^2} + \frac{\partial^2 u}{\partial z^2} \right) + f(x, y, z, t),
\]
где $c$ -- скорость распространения волн в данной среде, $f$ -- источник.

2) Процессы распространения тепла в однородном изотропном теле описываются уравнением теплопроводности
\[
\frac{\partial u}{\partial t} = a^2 \left( \frac{\partial^2 u}{\partial x^2} + \frac{\partial^2 u}{\partial y^2} + \frac{\partial^2 u}{\partial z^2} \right) + f(x, y, z, t),
\]
где $a^2$ – коэффициент температуропроводности данной среды.

3) При рассмотрении установившегося теплового состояния в однородном изотропном теле мы приходим к уравнению Пуассона:
\[
	\frac{\partial^2 u}{\partial x^2} + \frac{\partial^2 u}{\partial y^2} + \frac{\partial^2 u}{\partial z^2} = -\tilde{f}(x, y, z),
\]
где $\tilde{f}(x, y, z) = \frac{f(x, y, z)}{a^2}$

При отсутствии источников тепла внутри тела уравнение теплопроводности переходит в уравнение Лапласа:
\[
	\frac{\partial^2 u}{\partial x^2} + \frac{\partial^2 u}{\partial y^2} + \frac{\partial^2 u}{\partial z^2} = 0,
\]

Волновое уравнение (ВУ), уравнение теплопроводности (УТ), уравнение Пуассона (УП) и уравнение Лапласа (УЛ) часто называют основными уравнениями математической физики.

\bigskip

Каждое из основных уравнений математической физики имеет бесчисленное множество частных решений. При решении конкретной физической задачи, необходимо из всех этих решений выбрать то, которое удовлетворяет некоторым дополнительным условиям, вытекающим из её физического смысла. Итак, задачи математической физики состоят в отыскании решений уравнений в частных производных, удовлетворяющим некоторым дополнительным условиям. Такими дополнительными условиями чаще всего являются так называемые граничные (краевые) условия, т.е. условия, заданные на границе рассматриваемой среды, и начальные условия, относящегося к одному какому-нибудь моменту времени, с которого начинается изучение данного физического явления.

\bigskip

Обратим внимание на одно важное обстоятельство.

\textsl{Задача математической физики считается поставленной корректно (в смысле Адамара), если решение задачи, удовлетворяющее всем её условиям, существует, единственно и устойчиво} (непрерывно зависит от данных (параметров) задачи, т.е. малые изменения любого из данных задачи должно вызывать соответственно малое изменение решения). Требование устойчивости выражает физическую определенность поставленной задачи.

\bigskip

Дадим необходимые определения.

\begin{definition}[Дифференциальное уравнение в частных производных второго порядка (ДУЧП2П)]
Пусть $G$ --- область в $\mathbb{R}^n$, x=($x_1$, \ldots, $x_n$) --- точка из $G$. Дифференциальным уравнением в частных производных второго порядка (ДУЧП2П) с независимыми переменными $x_1, \ldots, x_n$ называется соотношение, связывающее значение неизвестной функции n переменных $u=u(x_1, \ldots, x_n)$, и её частных производных 1\textsuperscript{го} и 2\textsuperscript{го} порядков:
\[
\frac{\partial u}{\partial x_1}, \ldots, \frac{\partial u}{\partial x_n},
\frac{\partial^2 u}{\partial^2 x_1}, \ldots, \frac{\partial^2 u}{\partial x_i \partial x_j}, \ldots \frac{\partial^2 u}{\partial^2 x_n}
\] вида:
\[
F(x_1, \ldots, x_n, u, \frac{\partial u}{\partial x_1}, \ldots, \frac{\partial u}{\partial x_n},
\frac{\partial^2 u}{\partial^2 x_1}, \ldots, \frac{\partial^2 u}{\partial x_i \partial x_j}, \ldots \frac{\partial^2 u}{\partial^2 x_n})=0
\]
\end{definition}

Решением ДУЧП2П называется функция  класса , которая удовлетворяет данному уравнению.

Уравнение называется \textsl{линейным} относительно старших производных – производных второго порядка (ДУЧП2ПЛОСП), если оно имеет вид:

\[
\sum_{i,j=1}^{n} a_{ij}(x)\frac{\partial^2u}{\partial x_i \partial x_j} + \Phi(x, u, \operatorname{grad}(u)) = 0
\]
где коэффициенты при старших (вторых) производных $a_{ij}(x) \in C(G) (i, j = 1, \ldots,n)$ являются функциями только независимых переменных $x_1, \ldots, x_n$; 
а $\operatorname{grad}(u) = (\frac{\partial u}{\partial x_1}, \ldots, \frac{\partial u}{\partial x_n})$ --- градиент функции $u(x) = u(x_1, \ldots, x_n)$.

\bigskip

Коэффициенты при старших (вторых) производных образуют квадратную матрицу $A(x) = (a_{ij}(x))$  – матрицу старших коэффициентов:



\begin{note}
$\operatorname{grad}(u) = (\frac{\partial u}{\partial x_1}, \ldots, \frac{\partial u}{\partial x_n})$.
\end{note}
\begin{note}
Если $\Phi$ - линейная относительно $u(x)$, $\operatorname{grad}(u)$, то получаем линейное уравнение в частных производных 2\textsuperscript{го} порялка(ЛУЧП2П):
$$\sum_{i,j=1}^{n} a_{ij}(x)\frac{\partial^2u}{\partial x_i \partial x_j} + \sum_{i=1}^{n} b_{i}(x)\frac{\partial u}{\partial x_i} + c(x)u = f(x),$$
\[
  x \in G,\quad a_{ij}(x), b_i(x), c(x) \in C(G); \qquad\qquad (i,j = \overline{1,n})
\]
$$f(x) \in C(G)$$
\end{note}
Иногда пишут:
$L[u]=f(x)$, где
$$\displaystyle L= \sum_{i,j=1}^{n} a_{ij}(x)\frac{\partial^2}{\partial x_i \partial x_j} + \sum_{i=1}^{n} b_{i}(x)\frac{\partial }{\partial x_i} + c(x)$$
$\displaystyle L: C^2(G)\mapsto C(G)$ -- линейный дифференциальный оператор 2\textsuperscript{го} порядка(ЛДО2П).\par\medskip


\label{abbr:LDUChPSPPK}
\newcommand{\LDUChPSPPK}{\hyperref[abbr:LDUChPSPPK]{ЛДУЧП2ППК}}
Если же все коэффициенты линейного дифференциального уравнения частных производных 2\textsuperscript{го} порядка(ЛДУЧП2П): не зависят от $\vec x$, то получаем ЛДУЧП2П с постоянными коэффициентами (\LDUChPSPPK{}):

\[
\sum_{i,j=1}^{n} a_{ij}\frac{\partial^2u}{\partial x_i \partial x_j} + \sum_{i=1}^{n} b_{i}\frac{\partial u}{\partial x_i} + cu = f(x), \quad x \in G,
\]
где \(a_{ij}, b_i, c\) --- коэффициенты, \(i,j = \overline{1,n}\), \(f(x) \in C(G)\) --- правая часть. Это уравнение можно записать как \(L[u]=f\).\\
Здесь
$\displaystyle L= 
  \sum_{i,j=1}^{n} a_{ij}\frac{\partial}{\partial x_i} \frac{\partial}{\partial x_j}
  +\sum_{i=1}^{n} b_{i}\frac{\partial }{\partial x_i} + c$.\\ Первая сумма это \emph{квадратичная форма}, вторая -- \emph{линейная форма}.
\subsection{Линейная невырожденная замена переменных}
Осуществим в \LDUChPSPPK{} линейную невырожденную замену переменных:\\
$\displaystyle\vec x=\vec x(\vec \xi),$ где $\displaystyle\vec \xi=(\xi_1, \ldots, \xi_n),\quad \xi_k = \sum_{i=1}^{n}s_{ik}x_i,\ (k=\overline{1,n})$.\\
В матричной форме: $$\vec \xi = S^T \vec x$$
S --- матрица коэффициентов: S = ($s_{ik}$)$\Rightarrow$

\[
\begin{cases}
\xi_1 = s_{11}x_1 + \ldots + s_{n1}x_n \\
\ldots \\
\xi_i = s_{1i}x_1 + \ldots + s_{ni}x_n \\
\ldots \\
\xi_n = s_{1n}x_1 + \ldots + s_{nn}x_n
\end{cases}\Longrightarrow
\quad \det S \neq 0 \Longrightarrow \vec x=(S^{-1})^T\vec\xi
\]
Обозначим
\[
u(x(\xi))=\tilde{u}(\xi) \Rightarrow \tilde{u}(\xi)=u(x)\]
\[
\text{Т.к. } \frac{\partial \xi_k}{\partial x_i} = s_{ik} \text{, то} \frac{\partial u}{\partial x_i} = \sum_{k=1}^n\frac{\partial \tilde{u}}{\partial \xi_i}\frac{\partial \xi_k}{\partial x_i} = \sum_{k=1}^ns_{ik}\frac{\partial \tilde{u}}{\partial \xi_k}
\]


\begin{note}
Из $\displaystyle\frac{\partial }{\partial x_i} = \sum_{k=1}^n s_{ik}\frac{\partial}{\partial \xi_k}$ имеем:
\[\frac{\partial^2u}{\partial x_i \partial x_j}=\frac{\partial}{\partial x_i}(\frac{\partial u}{\partial x_j})=\sum_{l,k=1}^{n}\frac{\partial^2 \tilde{u}}{\partial\xi_k \partial\xi_l}s_{ik}s_{jl}\]
Подставляя полученные выражения для $\frac{\partial u}{\partial x_i}$ и $\frac{\partial^2 u}{\partial x_i \partial x_j}$ в \LDUChPSPPK{}, получим:

\[\xi_k=\sum_{i=1}^{n}s_{ik}x_i,\ \xi=S^Tx, x=(S^{-1})^T\xi\Rightarrow\]

\[
\sum_{k,l=1}^n \frac{\partial^2 \tilde{u}}{\partial \xi_k \partial \xi_l} \sum_{i,j=1}^n a_{ij}s_{ik}s_{jl} + \sum_{k=1}^n \frac{\partial \tilde{u}}{\partial \xi_k}\sum_{i=1}^nb_is_{ik}+c\tilde{u}=\tilde{f}
\]
Обозначим через 
\[\tilde{a}_{kl}=\sum_{i,j=1}^na_{ij}s_{ik}s_{jl}, \quad\tilde{b}_k=\sum_{i=1}^nb_is_{ik}\]
получаем:
$\displaystyle
\tilde{L}\tilde{u}=\sum_{k,l=1}^n\tilde{a}_{kl}\frac{\partial^2\tilde{u}}{\partial\xi_k\partial\xi_l}+\sum_{k=1}^n\tilde{b}_k\frac{\partial\tilde{u}}{\partial\xi_k}+c\tilde{u}=\tilde{f}(\xi),
$\\
где $\tilde{f}(\xi)=f(x(\xi))$\\
При этом матрица $A=(a_{ik})$ старших коэффициентов исходного ЛУЧП2ППК преобразуется по формулам\\
$
\displaystyle
  \tilde{A}=S^TAS 
$
   (из линейной алгебры)\\
Т.е. при невырожденной линейной замене переменных:
\[
\xi=S^Tx, detS\ne0,
\] матрица A преобразуется по формулам
\[
K=\sum_{i,j=1}^na_{ij}y_iy_j=y^TAy,\ y=(y_1, \ldots, y_n)^T
\]
\end{note}

\begin{note}
считаем A --- симметричной ($A^T=A$)
\end{note}
